\section{Procedimiento de implementación}

\subsection{Criterio aplicado en cada etapa}

Cada etapa se documentó con el mismo criterio: qué se hizo, por qué se eligió, qué resultado se esperaba, cómo se verificó y con qué objetivo se relaciona. La Tabla~\ref{tab:trazabilidad-etapas} resume esa trazabilidad antes de presentar el detalle y las evidencias.

\small
\begin{longtable}{p{0.13\linewidth}p{0.25\linewidth}p{0.22\linewidth}p{0.28\linewidth}}
\caption{Trazabilidad entre decisiones, resultados esperados, pruebas y objetivos.}\label{tab:trazabilidad-etapas}\\
\toprule
\textbf{Etapa} & \textbf{Qué se hizo y por qué} & \textbf{Resultado esperado} & \textbf{Verificación y relación con objetivos} \\
\midrule
\endfirsthead
\toprule
\textbf{Etapa} & \textbf{Qué se hizo y por qué} & \textbf{Resultado esperado} & \textbf{Verificación y relación con objetivos} \\
\midrule
\endhead
Hipervisor & Proxmox se instaló directamente en la torre para obtener un hipervisor Tipo 1 y evitar virtualización anidada. & Tres VM simultáneas sobre el hardware físico. & \texttt{pveversion}, recursos del host y panel de VM; cumple el objetivo de virtualización y la opción extra. \\
Sistemas & Ubuntu Server 26.04.1 se instaló sin GUI en web y DB; openSUSE se eligió como cliente gráfico distinto de Ubuntu. & Dos servidores de texto y una estación gráfica de administración. & \texttt{/etc/os-release}, \texttt{systemctl get-default} y KDE; cumple el requisito de SO y cliente. \\
Red y SSH & Se usó \texttt{vmbr1}, direccionamiento estático y SSH 22222 con llaves para administrar sin consola local. & Conectividad reproducible y acceso remoto cifrado a ambas VM. & Ping, \texttt{ss}, sesiones SSH desde \texttt{10.50.5.10} y puerto 22 cerrado; cumple gestión remota. \\
Base de datos & MySQL se separó en \texttt{lab5-db}, se ligó a \texttt{10.50.5.30:13306} y se limitó \texttt{wpapp} a la VM web. & WordPress accede a su DB; el cliente no puede hacerlo directamente. & Consulta positiva desde web, \texttt{SHOW GRANTS} y \texttt{nc} negativo desde openSUSE; cumple dos capas y mínimo privilegio. \\
Web y TLS & Apache, PHP y WordPress se instalaron en web; DNS-01 se eligió porque la IP del servicio es privada. & Sitio por nombre, redirección HTTP y certificado público válido. & \texttt{getent}, respuestas 301/200 y certificado Let's Encrypt; cumple resolución local y HTTPS. \\
SFTP y permisos & \texttt{wpadmin:www-data}, modos 750/640 y SFTP permiten administrar contenido sin conceder escritura general. & \texttt{wpadmin} modifica; usuarios no autorizados fallan; Apache conserva lectura. & FileZilla, \texttt{stat} y prueba negativa con \texttt{wpviewer}; cumple \texttt{chown}, \texttt{chmod} y restricción. \\
Firewall & UFW aplica denegación entrante y reglas por IP porque cada capa necesita flujos diferentes. & Solo permanecen accesibles los puertos y orígenes diseñados. & \texttt{ufw status verbose}, puertos en escucha y pruebas positivas/negativas; cumple seguridad de red. \\
\bottomrule
\end{longtable}
\normalsize

\subsection{Creación y preparación de las VM}

En Proxmox se crearon manualmente las VM 101, 102 y 103. Las dos VM servidoras recibieron Ubuntu Server 26.04.1 LTS y se fijó \texttt{multi-user.target}; el cliente recibió openSUSE Leap 16.0 y KDE. Los servidores se administran mediante cuentas sin privilegios directos y \texttt{sudo} solo cuando una tarea lo exige.

\evidencia{Ubuntu Server 26.04.1 LTS en la VM web}{03-servidores/E04-web-ubuntu-version.png}{La captura muestra \texttt{lab5-web}, la versión Ubuntu 26.04.1 LTS y \texttt{multi-user.target}. La ausencia de \texttt{graphical.target} satisface la restricción de servidor sin GUI.}

\evidencia{Ubuntu Server 26.04.1 LTS en la VM DB}{03-servidores/E05-db-ubuntu-version.png}{La misma comprobación en \texttt{lab5-db} confirma versión y objetivo de arranque de texto, independiente del servidor web.}

\subsection{Administración SSH}

En ambos servidores OpenSSH escucha en \texttt{22222/TCP}. La configuración deshabilita autenticación por contraseña y acceso remoto de \texttt{root}; las llaves ed25519 permanecen en \texttt{lab5-cliente}. La VM web admite \texttt{labadmin} y \texttt{wpadmin}; la VM DB admite la cuenta administrativa de base de datos. El puerto 22 se mantiene cerrado.

\begin{lstlisting}[caption={Controles principales de OpenSSH.}]
Port 22222
PermitRootLogin no
PasswordAuthentication no
PubkeyAuthentication yes
AllowUsers labadmin wpadmin
\end{lstlisting}

\evidencia{Acceso SSH a la VM web por 22222}{05-seguridad/E07-ssh-web-22222.png}{La sesión se origina en openSUSE y llega a \texttt{lab5-web} por el puerto alternativo. La comprobación de escucha permite verificar que 22/TCP no quedó expuesto.}

\evidencia{Acceso SSH a la VM DB por 22222}{05-seguridad/E08-ssh-db-22222.png}{La evidencia registra la sesión remota hacia \texttt{lab5-db}; la administración no depende de la consola de Proxmox.}

\subsection{MySQL en la segunda capa}

MySQL se configuró con \texttt{bind-address=10.50.5.30}, puerto \texttt{13306}, \texttt{skip-name-resolve=ON} y X Protocol desactivado. Así, 3306 y 33060 no quedan en escucha. La base \texttt{wordpress} utiliza el usuario \texttt{wpapp@10.50.5.20}; este usuario posee privilegios sobre \texttt{wordpress.*}, no privilegios globales.

\begin{lstlisting}[caption={Parámetros efectivos de escucha de MySQL.}]
[mysqld]
bind-address = 10.50.5.30
port = 13306
skip-name-resolve = ON
mysqlx = OFF
\end{lstlisting}

\evidencia{MySQL escucha únicamente en 10.50.5.30:13306}{04-wordpress-mysql/E09a-mysql-escucha-13306.png}{La salida de \texttt{ss} y las variables del servidor confirman la IP y el puerto alternativo. No existe escucha en 3306 ni en todas las interfaces.}

\evidencia{Conexión de WordPress hacia la VM DB}{04-wordpress-mysql/E09-mysql-remoto.png}{La consulta se ejecuta desde la VM web contra \texttt{lab5-db:13306}; demuestra el flujo positivo entre las dos capas.}

\evidencia{Privilegios limitados de wpapp}{04-wordpress-mysql/E09c-show-grants.png}{\texttt{SHOW GRANTS} muestra el alcance \texttt{wordpress.*} y el origen \texttt{10.50.5.20}; no se incluye contraseña ni permiso global.}

\subsection{WordPress y Apache}

WordPress se instaló en \texttt{/var/www/wordpress}. \texttt{wp-config.php} apunta a \texttt{10.50.5.30:13306} y fija \texttt{WP\_HOME} y \texttt{WP\_SITEURL} en \url{https://laboratorio5.com}. Apache usa un VirtualHost en 80 que redirige permanentemente a HTTPS y otro en 443 con el certificado de Let's Encrypt. Antes de cada recarga se ejecutó \texttt{apache2ctl configtest}.

El funcionamiento final no depende de una captura que muestre el panel administrativo o información sensible. Se demuestra mediante tres pruebas independientes incluidas en el informe: la consulta positiva de la VM web hacia MySQL, la respuesta HTTPS 200 obtenida desde openSUSE y el certificado válido para \texttt{laboratorio5.com}. En conjunto verifican Apache--PHP--WordPress--MySQL sin revelar contraseñas, llaves ni credenciales.

\resultado{WordPress respondió con código 200 por HTTPS, la VM web consultó la base remota \texttt{wordpress} en \texttt{10.50.5.30:13306} y el certificado presentado perteneció a \texttt{laboratorio5.com}.}
